\section{Introduction}\label{sec:intro}

Prediction markets aggregate dispersed beliefs into a single price
that, in equilibrium, behaves like a probability. The empirical
literature on these markets has historically focused on price-level
questions: forecast accuracy, calibration against realised outcomes,
the longshot bias, and the extent to which informed and uninformed
traders coexist on the same venue
\citep{wolfers-zitzewitz-2004,manski-2006,page-clemen-2013}.
Yet microstructure is what determines the trading cost of holding an
informational position, and the wider the cost, the smaller the
informational signal that survives to the price. A prediction market
with noisy microstructure produces noisier prices than the
headline-level aggregation literature implicitly assumes.
The microstructure of prediction markets, however, has remained
under-studied: limit-order-book depth profiles, spread decompositions,
and the behaviour of liquidity providers near resolution were largely
unobservable on the early venues, which used scoring rules
(IEM, the logarithmic market scoring rule of \citealt{hanson-2007})
or sparse parimutuel pools. Two recent
Polymarket-specific contributions sit close to the present paper but
ask different questions. \citet{rahman-alchami-clark-2025} survey
decentralised prediction-market microstructure as a methodological
domain, framing the open questions across venues without bringing
matching tick-level data. \citet{tsang-yang-2026} run a single-event
time-series microstructure study on the 2024 US presidential
election market, documenting episode-level patterns around Biden's
withdrawal, the September debate, and the October whale trades, and
reporting a roughly order-of-magnitude decline in Kyle's $\lambda$
as the market matured; the unit of analysis is one market over many
time slices, with $\lambda$ estimated on aggregated rather than
tick-level data. The present paper is the cross-sectional, tick-level
complement: 600 markets observed simultaneously over a single 28-day
scrape window, with microstructure measures computed on the full
event tape and a direct on-chain trade record. To our knowledge,
neither authoritative on-chain trade direction nor a continuous
off-chain order-book archive at tick resolution has been brought to
bear at the cross-sectional scale we use here.
The classical microstructure references behind our measurement
choices are \citet{ohara-1995} and \citet{hasbrouck-2007} for
order-book theory, \citet{foucault-pagano-roell-2013} for the
liquidity-provider angle, and
\citet{huang-stoll-1997,madhavan-richardson-roomans-1997} for spread
decomposition.

Polymarket changed this. Since 2021 the venue has run a
limit-order-book exchange on Polygon, settling in USDC against an
on-chain conditional-token contract. Our data covers 52 calendar days
of the public WebSocket feed (2026-02-21 to 2026-04-15) at tick
resolution: 30,287,264,368 events across 385,198 distinct markets.
A 28-day overlap window (2026-02-28 to 2026-03-27) is mirrored on
chain, where the CTF Exchange V1 contract logs 255 million
\texttt{OrderFilled} events whose payloads identify both
counterparties and the aggressor side. We use the off-chain feed for
quote dynamics and the on-chain record for trade direction, and study
both on a pre-registered, stratified 600-market panel.

The paper has two empirical contributions, ordered by weight for the
literature.

\textit{Trade-direction inference from Polymarket's public feed is
noise-dominated.} Six standard direction-dependent measures
(effective spread, realized spread, Kyle's $\lambda$, Amihud, Roll,
Abdi-Ranaldo) need an aggressor sign for each trade, and empirical
practice on equity venues sources that sign from a quote-driven feed
via Lee-Ready or its variants
\citep{lee-ready-1991,ellis-michaely-ohara-2000}. The Polymarket
public feed does not expose enough information to do this reliably:
the \texttt{change\_side} field marks which side of the book moved,
not which side initiated the trade. Sign agreement between
feed-inferred and on-chain trade directions sits at $\sim 59\%$
volume-weighted across the top-100 stratum and four disjoint 7-day
windows, just above the 50\% chance baseline. On the comparable
subset of the top-100 panel, the effective half-spread changes sign
on 67\% of markets in a first 7-day window when feed-inferred
direction is swapped for the on-chain record, and Kyle's $\lambda$
changes sign on 60\%; in a second non-overlapping window the rates
fall to 50\% and 43\% respectively. Both windows sit in or below
the chance band, well short of the $\sim 80\%$ Lee-Ready accuracy
documented on Nasdaq. Any Polymarket microstructure result that
depends on trade direction therefore needs to source it from
on-chain \texttt{OrderFilled} events. We release a replication
package that performs the join.

\textit{Eight cross-sectional stylized facts on the panel.} Quoted
half-spread shows the longshot premium familiar from the racetrack
literature; the L2 depth profile is closer to a uniform geometric
grid than to the top-of-book pattern usually assumed for prediction
markets; quote-update timing does not cluster on Polygon block
boundaries to a meaningful degree; maker-wallet concentration is
diverse on most markets with a thin concentrated tail; effective
spreads differ by category; archive-ingestion latency is tightly
distributed around 41.5\,ms with a multi-second $p_{99}$ tail; the
self-counterparty wash share has a 1\% median and a 22\% maximum
across the panel; and depth contracts in the cross-section as
markets approach resolution at a within-category log-log slope of
$0.55$. None of these eight facts requires the on-chain join, and
each is reported on the full 600-market pre-registered panel.

The two contributions are complementary. The eight stylized facts
characterise Polymarket's microstructure on its own terms. The
measurement result sets the conditions under which any
trade-direction-dependent statement about Polymarket can be trusted.

Section~\ref{sec:limits} carries the most counter-intuitive finding;
readers more interested in Polymarket as an empirical object will
care most about SF1 (longshot premium), SF2 (depth concentration),
and SF8 (depth decay near resolution).
